\clearpage
\section{Anhang: Pipeline-Durchlauf in drei Auflösungen}
\label{app:durchlauf}

Dieser Anhang zeigt die Zwischenstände eines vollständigen Pipeline-Durchlaufs
für drei Auflösungen desselben Testdatensatzes (5~FPS, OPENCV-Kameramodell).
Jede Zeile entspricht einer Zwischenstufe der Verarbeitung, jede Spalte einer
Auflösung: links \texttt{720p}, in der Mitte \texttt{qHD} und rechts
\texttt{low}. Die Reihenfolge folgt dem Datenfluss der Pipeline. Als Eingabe
dient das Videomaterial des Alurohr-Referenzdatensatzes. Über den Prompt
\texttt{pipe} wird das Zielobjekt promptbasiert segmentiert. Ab diesem
Schritt verzweigt der Pfad in die drei Spalten der Auflösungsvarianten.

% Die Querformatseiten erhalten eine eigene Satzvorlage. Dadurch werden die
% normalen, hochkant gesetzten fancyhdr-Kopf- und Fußzeilen unterdrückt.
\newcommand{\durchlaufheader}{%
  \vspace*{-59.5pt}%
  {\small
  \noindent Projektarbeit -- Scan-to-BIM-Pipeline\hfill Studienbereich Geo\par
  \vspace{2pt}\hrule height 0.4pt\vspace{6pt}}%
  \vspace{59.5pt}}
\newcommand{\durchlauffooter}{%
  \par\vspace*{\fill}%
  \begin{center}\small\thepage\end{center}%
  \vspace*{-41.2pt}}
\newcommand{\durchlaufbild}[1]{%
  \makebox[\hsize][c]{\includegraphics[width=0.95\hsize]{#1}}}
\newcommand{\durchlaufmasken}[1]{%
  \begin{minipage}[t]{\hsize}
  \centering
  \includegraphics[width=0.95\linewidth]{Screenshots/Pipline_Durchlauf/#1/review_ideal/frame_00000_middle_panel.png}\par\vspace{1pt}%
  \includegraphics[width=0.95\linewidth]{Screenshots/Pipline_Durchlauf/#1/review_ideal/frame_00080_middle_panel.png}\par\vspace{1pt}%
  \includegraphics[width=0.95\linewidth]{Screenshots/Pipline_Durchlauf/#1/review_ideal/frame_00160_middle_panel.png}\par\vspace{1pt}%
  \includegraphics[width=0.95\linewidth]{Screenshots/Pipline_Durchlauf/#1/review_ideal/frame_00239_middle_panel.png}%
  \end{minipage}}

\begin{landscape}
\thispagestyle{empty}
\durchlaufheader
\small
\renewcommand{\arraystretch}{1.05}
\setlength{\tabcolsep}{3pt}
\noindent
\begin{tabularx}{\linewidth}{@{}>{\raggedright\arraybackslash}p{2.4cm}>{\centering\arraybackslash}X>{\centering\arraybackslash}X>{\centering\arraybackslash}X@{}}
\toprule
Zwischenstufe & \textbf{\texttt{720p}} & \textbf{\texttt{qHD}} & \textbf{\texttt{low}}\\
\midrule
Masken (Prüfpanels) &
\durchlaufmasken{720p_Panels} &
\durchlaufmasken{qhd_panels} &
\durchlaufmasken{low_panels} \\
\bottomrule
\end{tabularx}
\par\vspace{3pt}
\noindent\parbox{\linewidth}{\footnotesize\textit{Maskenprüfung bei
\texttt{720p}, \texttt{qHD} und \texttt{low}. Je Auflösung sind die vier
Prüfzeitpunkte 00000, 00080, 00160 und 00239 untereinander dargestellt; jedes
Panel zeigt das Originalbild, die überlagerte \texttt{middle}-Maske und den
Objektausschnitt.}}
\durchlauffooter
\newpage

\thispagestyle{empty}
\durchlaufheader
\small
\renewcommand{\arraystretch}{1.05}
\setlength{\tabcolsep}{3pt}
\noindent
\begin{tabularx}{\linewidth}{@{}>{\raggedright\arraybackslash}p{2.4cm}>{\centering\arraybackslash}X>{\centering\arraybackslash}X>{\centering\arraybackslash}X@{}}
\toprule
Zwischenstufe & \textbf{\texttt{720p}} & \textbf{\texttt{qHD}} & \textbf{\texttt{low}}\\
\midrule
Punktwolken &
\durchlaufbild{Screenshots/Pipline_Durchlauf/Punktwolke_OpenCV_720p.png} &
\durchlaufbild{Screenshots/Pipline_Durchlauf/Punktwolke_OpenCV_qhd.png} &
\durchlaufbild{Screenshots/Pipline_Durchlauf/Punktwolke_OpenCV_low.png} \\
Punktwolken (Zoom) &
\durchlaufbild{Screenshots/Pipline_Durchlauf/Punktwolke_OpenCV_720p_zoom.png} &
\durchlaufbild{Screenshots/Pipline_Durchlauf/Punktwolke_OpenCV_qhd_zoom.png} &
\durchlaufbild{Screenshots/Pipline_Durchlauf/Punktwolke_OpenCV_low_zoom.png} \\
\bottomrule
\end{tabularx}
\par\vspace{3pt}
\noindent\parbox{\linewidth}{\footnotesize\textit{COLMAP-Punktwolken und vergrößerte Detailansichten
bei \texttt{720p}, \texttt{qHD} und \texttt{low}.}}
\durchlauffooter
\newpage

\thispagestyle{empty}
\durchlaufheader
\small
\renewcommand{\arraystretch}{1.05}
\setlength{\tabcolsep}{3pt}
\noindent
\begin{tabularx}{\linewidth}{@{}>{\raggedright\arraybackslash}p{2.4cm}>{\centering\arraybackslash}X>{\centering\arraybackslash}X>{\centering\arraybackslash}X@{}}
\toprule
Zwischenstufe & \textbf{\texttt{720p}} & \textbf{\texttt{qHD}} & \textbf{\texttt{low}}\\
\midrule
Splat (Opazität 9999) &
\durchlaufbild{Screenshots/Pipline_Durchlauf/Splat_OpenCV_720p_9999opac.png} &
\durchlaufbild{Screenshots/Pipline_Durchlauf/Splat_OpenCV_qhd_9999opac.png} &
\durchlaufbild{Screenshots/Pipline_Durchlauf/Splat_OpenCV_low_9999opac.png} \\
Mesh &
\durchlaufbild{Screenshots/Pipline_Durchlauf/Splats_to_Mesh_720p.png} &
\durchlaufbild{Screenshots/Pipline_Durchlauf/Splats_to_Mesh_qhd.png} &
\durchlaufbild{Screenshots/Pipline_Durchlauf/Splats_to_Mesh_low.png} \\
\bottomrule
\end{tabularx}
\par\vspace{3pt}
\noindent\parbox{\linewidth}{\footnotesize\textit{Splat- und Mesh-Zwischenstände bei
\texttt{720p}, \texttt{qHD} und \texttt{low}.}}
\durchlauffooter
\newpage

\thispagestyle{empty}
\durchlaufheader
\small
\renewcommand{\arraystretch}{1.05}
\setlength{\tabcolsep}{3pt}
\noindent
\begin{tabularx}{\linewidth}{@{}>{\raggedright\arraybackslash}p{2.4cm}>{\centering\arraybackslash}X>{\centering\arraybackslash}X>{\centering\arraybackslash}X@{}}
\toprule
Zwischenstufe & \textbf{\texttt{720p}} & \textbf{\texttt{qHD}} & \textbf{\texttt{low}}\\
\midrule
Mesh (ohne Nachbearbeitung) &
\durchlaufbild{Screenshots/Pipline_Durchlauf/Splats_to_Mesh_720p_ugly.png} &
\durchlaufbild{Screenshots/Pipline_Durchlauf/Splats_to_Mesh_qhd_ugly.png} &
\durchlaufbild{Screenshots/Pipline_Durchlauf/Splats_to_Mesh_low_ugly.png} \\
Centerline &
\durchlaufbild{Screenshots/Pipline_Durchlauf/Centerline_720p.png} &
\durchlaufbild{Screenshots/Pipline_Durchlauf/Centerline_qhd.png} &
\durchlaufbild{Screenshots/Pipline_Durchlauf/Centerline_low.png} \\
\bottomrule
\end{tabularx}
\par\vspace{3pt}
\noindent\parbox{\linewidth}{\footnotesize\textit{Unbearbeiteter Mesh-Rohstand und extrahiertes
Linienelement (Centerline) als Endstufen der Pipeline bei \texttt{720p},
\texttt{qHD} und \texttt{low}.}}
\durchlauffooter
\end{landscape}
